\scanpage{380}
\mspage{104}
% archive sheet number: 379
bon usage, et en déduire le \deleted{\unsure{XV}} Th.~10.1 (XV, p.~134-136). Mieux que l'existence des $f_*$, l'ensemble d'axiomes \unsure{étonnamment} faibles que je viens d'expliciter \unsure{suffit} à assurer que $\mathrm{HOT}_W$ est un \uline{dérivateur} sur $\mathrm{Cat}$ \deleted{\unsure{tout|tous}} \unsure{entier|entière}, au plein sens du terme, i.e. qu'il satisfait à $\mathrm{Der}(1,2,3,\deleted{\unsure{4^*}}4,\deleted{\unsure{5^*}}5,6,4^*,5^*,6^*)$.

Ceci est obtenu en remplaçant \inserted{(grâce à XV 8.4)} $\mathrm{Cat}$ par $A^{\wedge}$, où $A$ est une catégorie-test convenable, \inserted{$W$ par $W_A\subset\mathrm{Fl}\,A^{\wedge}$,} et en prouvant que \deleted{la} \deleted{\unsure{dér}} $\unsure{W_A|W/A}$ s'insère dans un triple de Quillen clos $(W,C,F)$, avec $C=\mathrm{Mono}(A^{\wedge})$, ce qui détermine $F$ comme $(C\cap W_A)_*$. On applique alors $\underline{\overline{\mathrm{XV}}}$ 7.4. à cette situation (résultats résumés dans XIX th.~2.3 (p.~19), 2.6 (p.~25) 2.9 (p.~28).
On retrouve l'existence des $f_!$ (avec un mode de calcul bien différent du \unsure{calcul|celui} dans $\mathrm{Cat}$ !), et en plus on trouve l'existence des $f_*$ (plus l'exactitude des carrés
